\documentclass[10pt,a4paper]{amsart}
\usepackage[margin=19mm]{geometry}
\usepackage{amsmath,amssymb,mathtools}
\usepackage[scr=rsfs,cal=euler]{mathalfa}
\usepackage{xfrac}
\usepackage[unicode,hidelinks,bookmarks=false]{hyperref}
\newcommand{\ie}{i.e.\ }
\newcommand{\cf}{cf.\ }
\newcommand{\resp}{respectively\ }
\newcommand{\Subsection}{\subsection}
\providecommand{\nobreakdash}{\nobreak\hbox{-}}
\setlength{\emergencystretch}{2em}
\multlinegap0pt
\usepackage{amssymb}
\usepackage[shortlabels]{enumitem}
\newtheorem{theorem}{Theorem}
\newcommand{\B}{\mathbb{B}}
\newcommand{\C}{\mathcal{C}}
\newcommand{\CC}{\mathbb{C}}
\newcommand{\LL}{\mathbb{L}}
\newcommand{\M}{\mathbb{M}}
\newcommand{\N}{\mathbb{N}}
\newcommand{\inn}{~ \hat{\in}~ }
\title{Complex circles of partition and the expansion principles}
\author{Berndt Gensel, Theophilus Agama}
\date{Source: arXiv:2304.13371v6 (CC BY 4.0)}
\begin{document}
\maketitle

\noindent\emph{Source and licence.} Complex circles of partition and the expansion principles. Version arXiv:2304.13371v6 (CC BY 4.0), distributed by arXiv under the Creative Commons Attribution 4.0 International licence, http://creativecommons.org/licenses/by/4.0/, as recorded in the arXiv licence metadata for this exact version and shown on the abstract page https://arxiv.org/abs/2304.13371v6. Full licence notice: \texttt{licenses/arxiv-2304.13371-CC-BY-4.0.txt}.\par\smallskip

\noindent\textit{Editorial scope.} Counted unit: the article's theorem T-Limits (source0.tex lines 821--835). The article's complete original proof of it is delivered with it (source0.tex lines 837--877, one \texttt{proof} environment of 41 lines). No mathematics is written, repaired or re-derived: the statement and the proof are the article's own lines, in the article's own notation, and no retained line is substituted or re-flowed.  Retained uncounted as the source-local context the proof needs: lines 710--725.  Nothing was omitted from any retained span, so the package carries no omission record.  No retained line contains a citation, so the wrapper carries no bibliography and no bibliography entry.  No retained line contains a figure environment, an image-inclusion command or drawing code, so no image is delivered and the package declares no asset dependency.  Editorial formatting only, in addition: an AMS \texttt{amsart} preamble that restates the article's own declarations and macros used by the retained lines, namely the environments \texttt{theorem} on separate counters, together with the standard packages those lines need.  The retained environments print in this document as Theorem 1 in order of appearance, whereas the article prints the same items with the numbering of its own full document.  The complete native source of the article as distributed in the arXiv e-print for this exact version is bundled unchanged as \texttt{sources/140609/source0.tex}.
\begin{theorem}[The expansion principles]\label{T-Expansion}
Let $\B\subset\M\subseteq\N$ and $n,t,s\in\mathbb{Q}\subseteq\N$. Furthermore, let $\C^o(n,\CC_{\M})$ and $\C^o(n+t,\CC_{\M})$
be two nonempty cCoPs with the axes 
\[
\LL_{[z],[n-z]}\inn\C^o(n,\CC_{\M})\mbox{ and }\LL_{[w],[n+t-w]} \inn \C^o(n+t,\CC_{\M}). 
\]
If 
\[
w,n-z\in\CC_{\B}\mbox{ and }\Re(w)=\Re(z)+s, 
\]
then the axis $\LL_{[u],[n+s-u]}\inn \C^o(n+s,\CC_{\B})$ exists with
\[
\Re(u)=\Re(w) \mbox{ and }\Re(n+s-u)=\Re(n-z)
\]
and the imaginary parts corresponding to the \emph{circle condition}. Thus, $\C^o(n+s,\CC_{\B})$ is a nonempty cCoP.
\end{theorem}

\begin{theorem}\label{T-Limits}
Let $\M\subseteq\N$ and $\C^o(n,\CC_{\M})$ and $\C^o(n+t,\CC_{\M})$ be nonempty cCoP with positive integers $n,t\in\mathbb{Q}\subseteq\N$. If there is an axis $\LL_{[z_1],[z_2]}$ of $\C^o(n,\CC_{\M})$, then there are--depending on these axis--a lower (A:) and an upper (B:) limiting axis $\LL_{[u_1],[u_2]}$ (resp.  $\LL_{[v_1],[v_2]}$) of the cCoP $\C^o(n+t,\CC_{\M})$ for their axes $\LL_{[w_1],[w_2]}$ with $\Re(w_1)>\Re(z_1)$ and
\begin{align}
\Re(u_1)&=\Re(z_1) \mbox{ and } \Re(v_2)=\Re(z_2)\label{E-z1u1v2z2}
\end{align}
such that
\begin{align}
A: \Re(u_1)<\Re(w_1)<\Re(v_1)\text{ and }B: \Re(v_2)<\Re(w_2)<\Re(u_2)\label{E-ForSqueeze}
\end{align}
for the Squeeze Principle and
\begin{align}
A: \Re(v_1)<\Re(w_1)<\frac{n+t}{2}\text{ and }B: \frac{n+t}{2}<\Re(w_2)<\Re(v_2)\label{E-ForForecast}
\end{align}
for the Forecast Principle.
\end{theorem}

\begin{proof}
Due to $w_1,z_2\in\CC_B$ and $\LL_{[w_1],[z_2]}\inn \C^o(n+s,\CC_B)$, we deduce $\Re(w_1)+\Re(z_2)=n+s$. Furthermore, we use the requirement of Theorem \ref{T-Expansion}: $\Re(w_1)=\Re(z_1)+s$, so that with \eqref{E-z1u1v2z2} and (\ref{E-ForSqueeze}.A), we get
\begin{align*}
\Re(u_1)&<\Re(w_1)<\Re(v_1)\\
\Re(z_1)&<\Re(w_1)<n+t-\Re(v_2)\\
\Re(z_1)&<\Re(w_1)<n+t-\Re(z_2)\mid +\Re(z_2)\\
\Re(z_1)+\Re(z_2)&<\Re(w_1)+\Re(z_2)<n+t\\
n&<n+s<n+t.
\end{align*} 
We deduce $0<s<t$, which implies the \emph{squeeze principle}.\\

Also, by \eqref{E-z1u1v2z2} and (\ref{E-ForSqueeze}.B), we get
\begin{align*}
\Re(v_2)&<\Re(w_2)<\Re(u_2)\\
\Re(z_2)&<n+t-\Re(w_1)<n+t-\Re(u_1)\mid +\Re(w_1)\\
\Re(z_2)+\Re(w_1)&<n+t<n+t-\Re(z_1)+\Re(w_1)\\
n+s&<n+t<n+t+s\text{, since }\Re(w_1)-\Re(z_1)=s\\
s&<t<t+s.
\end{align*}
We deduce $t>s$, which implies the \emph{squeeze principle}.\\

Again, by \eqref{E-z1u1v2z2} and (\ref{E-ForForecast}.A), we get
\begin{align*}
\Re(v_1)&<\Re(w_1)<\frac{n+t}{2}\\
n+t-\Re(v_2)&<\Re(w_1)<\frac{n+t}{2}\\
n+t-\Re(z_2)&<\Re(w_1)<\frac{n+t}{2}\mid +\Re(z_2) \\
n+t&<\Re(w_1)+\Re(z_2)<\frac{n+t}{2}+\Re(z_2)\\
n+t&<\Re(w_1)+\Re(z_2)<\frac{n+t}{2}+n\\
n+t&<n+s<\frac{n+t}{2}+n.
\end{align*}
We deduce $t<s<\frac{n+t}{2}$, which implies the \emph{forecast principle}.\\

Similarly, by \eqref{E-z1u1v2z2} and (\ref{E-ForForecast}.B), we get
\begin{align*}
\frac{n+t}{2}&<\Re(w_2)<\Re(v_2)\\
\frac{n+t}{2}&<n+t-\Re(w_1)<\Re(v_2)=\Re(z_2)\mid +\Re(w_1)\\
\frac{n+t}{2}+\Re(w_1)&<n+t<\Re(w_1)+\Re(z_2)\\
\frac{n+t}{2}&<n+t<n+s.
\end{align*}
We deduce $t<s$, which implies the \emph{forecast principle}. This completes the proof.
\end{proof}

\end{document}
